\section{Контроль якості й фокус цього продовження}
\sourceblock{Solà, \texttt{1711.02508}; micro-Lie, \texttt{1812.01537}; локальний контекст FetchQuest.}

Цей блок є не черговим оглядом, а щільним перекладним ядром для найбільш корисної частини бібліотеки: фільтри стану, геометрія груп Лі, VIO/SLAM-нев'язки, практичні формули Калмана і міст до чисельних методів. Новий пошук літератури тут свідомо не розширювався. Поточний корпус уже має пріоритетні сирцеві пакети: Solà ESKF, micro-Lie, Zuazua, PySDR, SDR survey, Correll robotics, Labbe Kalman, FDM/FEM та Open Optimization. Наступний виграш дає не ширина, а чисте машинозчитуване TeX-перекладення формул і навколишньої технічної прози.

Практичний критерій готовності модуля:
\begin{enumerate}
\item формули, позначення, посилання і змінні не зруйновано перекладом;
\item кожен блок має явний зв'язок з реалізацією: стан, шум, нев'язка, матриця Якобі, коваріація, перевірка одиниць;
\item український текст пояснює, що саме інженер має не переплутати: локальна/глобальна похибка, ліве/праве збурення, дискретизація шуму, порядок множення кватерніонів;
\item PDF компілюється XeLaTeX без зниклих кириличних гліфів, кліпінгу заголовків і переповнених таблиць.
\end{enumerate}

Поточний фокус: закрити найважчі вузли для подальших агентів. Це означає: Solà ESKF як еталон для IMU/VIO; micro-Lie як еталон для $\SO(3)$, $\SE(3)$ та фактор-графів; бібліотека нев'язок VIO/SLAM як міст до сучасних статей; Labbe Kalman як міст від інтуїтивного фільтрування до формальної ESKF-нотації.
